% ============================================================================
%  fig_conexiones.tex - Circuito de conexiones del prototipo F6 (nivel de pines)
%  Pong en FPGA - Electrónica Digital 1 - UNAL
%  Juan David Ortiz Muñoz, David Eduardo Chivata Arias, Daniel Felipe Castro González
%  Datos tomados de las tablas de conexión de la memoria de F6.
%  Compilar con: pdflatex fig_conexiones.tex
% ============================================================================
\documentclass[border=6pt]{standalone}
\usepackage[utf8]{inputenc}
\usepackage[T1]{fontenc}
\usepackage{lmodern}
\usepackage{newunicodechar}
\newunicodechar{Ω}{\ensuremath{\Omega}}
\usepackage[siunitx, american]{circuitikz}
\usetikzlibrary{arrows.meta,calc}
\begin{document}
\begin{circuitikz}[
    font=\sffamily\scriptsize,
    ic/.style={draw, line width=0.8pt, fill=yellow!10},
    mod/.style={draw, line width=0.8pt, fill=red!8},
    fpga/.style={draw, line width=1pt, fill=blue!5},
    pl/.style={font=\sffamily\tiny},          % número de pin dentro de la caja
    w/.style={line width=0.55pt},
    nota/.style={font=\sffamily\tiny, align=center, text=black!75},
]

% =========================================================== FPGA
\draw[fpga] (0, 1.2) rectangle (3.4, 15.4);
\node[font=\sffamily\small\bfseries, align=center] at (1.7, 15.0) {FPGA};
\node[align=center, font=\sffamily\tiny] at (1.7, 14.45) {EP4CE10E22C8N\\A-C4E6E10 · 3,3 V};

% Salidas (lado derecho): pin, señal, y
\foreach \p/\s/\y in {28/DIN/14.0, 31/CLK\_LED/13.4, 33/LOAD/12.8, 49/HCT\_EN/11.0,
                      101/J1\_BCD0/8.6, 103/J1\_BCD1/8.0, 104/J1\_BCD2/7.4, 105/J1\_BCD3/6.8,
                      106/J2\_BCD0/5.0, 110/J2\_BCD1/4.4, 111/J2\_BCD2/3.8, 112/J2\_BCD3/3.2}{
  \node[pl, anchor=east] at (3.35, \y) {\s\ \textbf{\p}};
}
\node[nota] at (1.7, 2.1) {BCD por el conector J2 (LCD),\\contactos 7–14};

% Entradas (lado izquierdo)
\foreach \p/\s/\y in {30/J1\_UP\_N/10.4, 32/J1\_DOWN\_N/9.9, 34/J1\_DETECT\_N/9.4,
                      39/J2\_UP\_N/8.6, 38/J2\_DOWN\_N/8.1, 42/J2\_DETECT\_N/7.6,
                      43/START\_N/6.6, 46/PAUSE\_N/6.1, 44/RESET\_N/5.6, 23/CLK\_SYS/4.4}{
  \node[pl, anchor=west] at (0.05, \y) {\textbf{\p}\ \s};
  \draw[w, ->] (-1.2, \y) -- (0, \y);
}
\draw[w] (-1.3, 9.2) -- (-1.4, 9.2) -- (-1.4, 10.6) -- (-1.3, 10.6);
\node[anchor=east, align=right] at (-1.5, 9.9) {Mando J1};
\draw[w] (-1.3, 7.4) -- (-1.4, 7.4) -- (-1.4, 8.8) -- (-1.3, 8.8);
\node[anchor=east, align=right] at (-1.5, 8.1) {Mando J2};
\draw[w] (-1.3, 5.4) -- (-1.4, 5.4) -- (-1.4, 6.8) -- (-1.3, 6.8);
\node[anchor=east, align=right] at (-1.5, 6.1) {Pulsadores\\a GND};
\node[anchor=east, align=right] at (-1.5, 4.4) {Oscilador\\50 MHz};
\node[nota, anchor=north] at (-1.6, 5.0) {};
\node[nota, align=left, anchor=north west] at (-3.3, 3.6) {Entradas \_N activas en bajo:\\reposo 3,3 V, pulsado GND.};

% =========================================================== SN74HCT245N #1 (matrices)
\draw[ic] (6.6, 12.3) rectangle (9.2, 14.6);
\node[font=\sffamily\scriptsize\bfseries] at (7.9, 14.3) {SN74HCT245N \#1};
\foreach \a/\b/\pa/\pb/\y in {A1/B1/2/18/14.0, A2/B2/3/17/13.4, A3/B3/4/16/12.8}{
  \node[pl, anchor=west] at (6.65, \y) {\pa\ \a};
  \node[pl, anchor=east] at (9.15, \y) {\b\ \pb};
  \draw[w] (3.4, \y) -- (6.6, \y);
}
\node[nota, anchor=south] at (7.9, 14.65) {VCC 20 y DIR 1 a +5 V · GND 10 · 100 nF};
\node[pl] at (7.9, 12.5) {19 /OE};

% =========================================================== Matrices FC-16
\draw[mod] (12.0, 12.3) rectangle (14.6, 14.6);
\node[font=\sffamily\footnotesize\bfseries, align=center] at (13.3, 14.35) {FC-16 \#1};
\draw[mod] (16.2, 12.3) rectangle (18.8, 14.6);
\node[font=\sffamily\footnotesize\bfseries, align=center] at (17.5, 14.35) {FC-16 \#2};
\foreach \s/\y in {DIN/14.0, CLK/13.4, CS/12.8}{
  \node[pl, anchor=west] at (12.05, \y) {\s};
}
\foreach \s/\so/\y in {DIN/DOUT/14.0, CLK/CLK/13.4, CS/CS/12.8}{
  \node[pl, anchor=east] at (14.55, \y) {\so};
  \node[pl, anchor=west] at (16.25, \y) {\s};
  \draw[w] (14.6, \y) -- (16.2, \y);
}
\draw[w] (9.2, 14.0) -- (12.0, 14.0);
\draw[w] (9.2, 13.4) -- (12.0, 13.4);
\draw[w] (9.2, 12.8) -- (12.0, 12.8);
\node[nota, anchor=south] at (10.6, 14.05) {10 kΩ a GND};
\node[nota, anchor=south] at (10.6, 13.45) {10 kΩ a GND};
\node[nota, anchor=south] at (10.6, 12.85) {10 kΩ a +5 V};
\node[pl, align=center] at (13.3, 12.65) {VCC +5 V · GND};
\node[pl, align=center] at (17.5, 12.65) {VCC +5 V · GND};
\node[nota, anchor=north] at (15.4, 12.2) {GC7219C (compatible MAX7219), cascada: tablero 8 × 16};

% =========================================================== Habilitación: 2 × 2N3904
% Rama #1
\draw (7.1, 11.0) node[npn, anchor=B] (q1) {};
\draw[w] (q1.C) -- (q1.C |- 0,12.3);
\draw (q1.E) node[ground] {};
\node[anchor=west, font=\sffamily\tiny, align=left] at ($(q1.C)+(0.1,-0.4)$) {Q1 2N3904\\/OE \#1 con 10 kΩ a +5 V};
% Rama #2
\draw (7.1, 1.6) node[npn, anchor=B] (q2) {};
\draw[w] (q2.C) -- (q2.C |- 0,2.6);
\draw (q2.E) node[ground] {};
\node[anchor=west, font=\sffamily\tiny, align=left] at ($(q2.C)+(0.1,-0.4)$) {Q2 2N3904\\/OE \#2 con 10 kΩ a +5 V};
% Base: HCT_EN con 10 kΩ en serie y 100 kΩ a GND
\draw[w] (3.4, 11.0) -- (4.4, 11.0);
\draw (4.4, 11.0) to[R, l_=\tiny 10 kΩ] (q1.B);
\fill (4.4, 11.0) circle (1.4pt);
\draw[w] (4.4, 11.0) -- (4.4, 9.8) -- (5.2, 9.8) -- (5.2, 1.6);
\draw (5.2, 1.6) to[R, l=\tiny 10 kΩ] (q2.B);
\node[nota, anchor=south] at (5.6, 11.25) {100 kΩ a GND en la base};
\node[nota, anchor=north] at (6.1, 1.2) {100 kΩ a GND en la base};

% =========================================================== SN74HCT245N #2 (marcador)
\draw[ic] (6.6, 2.6) rectangle (9.2, 9.4);
\node[font=\sffamily\scriptsize\bfseries] at (7.9, 9.1) {SN74HCT245N \#2};
\foreach \a/\b/\pa/\pb/\y in {A1/B1/2/18/8.6, A2/B2/3/17/8.0, A3/B3/4/16/7.4, A4/B4/5/15/6.8,
                              A5/B5/6/14/5.0, A6/B6/7/13/4.4, A7/B7/8/12/3.8, A8/B8/9/11/3.2}{
  \node[pl, anchor=west] at (6.65, \y) {\pa\ \a};
  \node[pl, anchor=east] at (9.15, \y) {\b\ \pb};
  \draw[w] (3.4, \y) -- (6.6, \y);
}
\node[pl] at (7.9, 2.8) {19 /OE};
\node[pl, align=center] at (7.9, 5.9) {VCC 20, DIR 1\\a +5 V · GND 10};

% =========================================================== CD4511BE y displays
\newcommand{\marcador}[3]{% #1 = nombre, #2 = y de la entrada A, #3 = y superior de la caja
  % decodificador
  \draw[ic] (11.4, #3-2.8) rectangle (14.0, #3);
  \node[font=\sffamily\scriptsize\bfseries, align=center] at (12.6, #3-1.4) {CD4511BE\\#1};
  \foreach \e/\p [count=\i from 0] in {A/7, B/1, C/2, D/6}{
    \pgfmathsetmacro{\yy}{#2-0.6*\i}
    \node[pl, anchor=west] at (11.45, \yy) {\p\ \e};
    \draw[w] (9.2, \yy) -- (11.4, \yy);
  }
  % salidas, resistencias y display
  \draw[mod] (17.0, #3-2.8) rectangle (18.8, #3);
  \node[font=\sffamily\footnotesize\bfseries, align=center] at (18.1, #3-1.4) {5161AS\\#1};
  \foreach \s/\po/\pd [count=\i from 0] in {a/13/7, b/12/6, c/11/4, d/10/2, e/9/1, f/15/9, g/14/10}{
    \pgfmathsetmacro{\yy}{#3-0.25-0.35*\i}
    \node[pl, anchor=east] at (13.95, \yy) {\s\ \po};
    \node[pl, anchor=west] at (17.05, \yy) {\pd\ \s};
    \draw[w] (14.0, \yy) -- (14.6, \yy);
    \draw[w] (15.8, \yy) -- (17.0, \yy);
    \draw[w] (14.6, \yy) -- ++(0.15, 0.08) -- ++(0.15,-0.16) -- ++(0.15, 0.16) -- ++(0.15,-0.16) -- ++(0.15, 0.16) -- ++(0.15,-0.16) -- ++(0.15, 0.08) -- (15.8, \yy);
  }
  \node[nota, anchor=south] at (15.2, #3-0.05) {7 × 470 Ω};
  \node[pl, align=center] at (18.4, #3-2.55) {3, 8\\GND};
}
\marcador{J1}{8.6}{9.4}
\marcador{J2}{5.0}{5.8}
\node[nota, align=left, anchor=north west] at (11.3, 2.85) {CD4511BE: VDD 16 a +5 V, VSS 8 a GND,\\LT 3 y BL 4 a +5 V, LE 5 a GND;\\entradas A–D con 10 kΩ a GND (marcador 0–0 sin FPGA).\\5161AS de cátodo común.};

% =========================================================== Alimentación
\node[draw, fill=red!6, align=center, font=\sffamily\scriptsize] at (17.5, 1.4)
  {Fuente externa 5 V / 3 A\\GND común con la FPGA};
\node[draw, fill=red!6, align=center, font=\sffamily\scriptsize] at (1.7, 0.4)
  {Mini-USB (3,3 V internos)};
\end{circuitikz}
\end{document}
